\section{Degree-three seed generation by the quadratic field}
\label{sec:cubic_generation}

The degree-three seed set is generated from the complete quadratic dictionary:
\[
\seedM_3^{\mathcalE}
=
\bigcup_{q\in\mathcal M_2^{\mathcalE}}
\Mon_3(\LieQ q).
\]
Only the quadratic transmission field is used in this step. The linear field
has already been assigned a different role: it closes the homogeneous block
after the seeds have been created.

\subsection{Two general differentiation formulas}

For an infection observable $S_iI_j$, the quadratic field acts only through the
susceptible factor, because $f_{\mathrm Q}$ has zero infectious component:
\[
\LieQ(S_iI_j)
=
-\beta\sum_{\ell\in\Min(i)}
 a_{i\ell}S_iI_\ell I_j.
\]
This creates monomials of family $SII$.

For an exposure observable $S_iE_j$, the product rule has two contributions:
\[
\LieQ(S_iE_j)
=
(\dot S_i|_{f_{\mathrm Q}})E_j
+
S_i(\dot E_j|_{f_{\mathrm Q}}).
\]
Therefore
\[
\LieQ(S_iE_j)
=
-\beta\sum_{\ell\in\Min(i)}
 a_{i\ell}S_iI_\ell E_j
+
\beta\sum_{\ell\in\Min(j)}
 a_{j\ell}S_iS_jI_\ell.
\]
The first sum creates $SIE$ monomials. The second creates $SSI$ monomials.

\subsection{How the terms arise: representative calculations}

Consider the diagonal observable $S_1I_1$. Since city 1 receives infection
from cities 1 and 2,
\[
\begin{aligned}
\LieQ(S_1I_1)
&=(\dot S_1|_{f_{\mathrm Q}})I_1\\
&=-\beta S_1(a_{11}I_1+a_{12}I_2)I_1\\
&=-\beta a_{11}S_1I_1^2
  -\beta a_{12}S_1I_1I_2.
\end{aligned}
\]
The repeated factor $I_1^2$ is retained. It is a genuine commutative monomial,
not a duplicate to be removed.

For the off-diagonal observable $S_1I_2$,
\[
\LieQ(S_1I_2)
=-\beta a_{11}S_1I_1I_2
 -\beta a_{12}S_1I_2^2.
\]
The monomial $S_1I_1I_2$ also appeared in the preceding derivative. It will be
represented by one cubic dictionary coordinate, but it will occur in two
different rows of the degree-two-to-degree-three coupling matrix.

Now consider the diagonal exposure observable $S_1E_1$:
\[
\begin{aligned}
\LieQ(S_1E_1)
={}&-\beta S_1(a_{11}I_1+a_{12}I_2)E_1\\
&+\beta S_1^2(a_{11}I_1+a_{12}I_2)\\
={}&-\beta a_{11}S_1E_1I_1
   -\beta a_{12}S_1E_1I_2\\
&+\beta a_{11}S_1^2I_1
   +\beta a_{12}S_1^2I_2.
\end{aligned}
\]
The negative contribution comes from differentiating the susceptible factor;
the positive contribution comes from differentiating the exposed factor.

Finally, for the off-diagonal observable $S_1E_2$,
\[
\begin{aligned}
\LieQ(S_1E_2)
={}&-\beta S_1(a_{11}I_1+a_{12}I_2)E_2\\
&+\beta S_1S_2(a_{21}I_1+a_{22}I_2+a_{23}I_3).
\end{aligned}
\]
Thus it generates two $SIE$ terms and three $SSI$ terms. The second group
contains information from the incoming sources of city 2 because the exposed
factor belongs to city 2.

\subsection{Complete generation table}

The following table lists the distinct cubic terms generated in each scalar
quadratic-observable equation. Coefficients are omitted from the table only to
make dictionary membership easier to read; the complete coefficient-bearing
formulas are given in \cref{app:full_derivatives}.

% Generated by scripts/generate_example_artifacts.py
\begin{longtable}{@{}>{\raggedright\arraybackslash}p{0.17\textwidth}>{\raggedright\arraybackslash}p{0.76\textwidth}@{}}
\caption{Degree-three monomials generated by $f_{\mathrm Q}$ from each quadratic coordinate.}
\label{tab:cubic_generation}\\
\toprule
Quadratic coordinate & Distinct cubic monomials in $\LieQ q$\\
\midrule
\endfirsthead
\toprule
Quadratic coordinate & Distinct cubic monomials in $\LieQ q$\\
\midrule
\endhead
$S_1I_1$ & $S_1I_1^{2}$,\ $S_1I_1I_2$ \\
$S_1I_2$ & $S_1I_1I_2$,\ $S_1I_2^{2}$ \\
$S_2I_1$ & $S_2I_1^{2}$,\ $S_2I_1I_2$,\ $S_2I_1I_3$ \\
$S_2I_2$ & $S_2I_1I_2$,\ $S_2I_2^{2}$,\ $S_2I_2I_3$ \\
$S_2I_3$ & $S_2I_1I_3$,\ $S_2I_2I_3$,\ $S_2I_3^{2}$ \\
$S_3I_2$ & $S_3I_2^{2}$,\ $S_3I_2I_3$ \\
$S_3I_3$ & $S_3I_2I_3$,\ $S_3I_3^{2}$ \\
$S_1E_1$ & $S_1E_1I_1$,\ $S_1E_1I_2$,\ $S_1^{2}I_1$,\ $S_1^{2}I_2$ \\
$S_1E_2$ & $S_1E_2I_1$,\ $S_1E_2I_2$,\ $S_1S_2I_1$,\ $S_1S_2I_2$,\ $S_1S_2I_3$ \\
$S_2E_1$ & $S_2E_1I_1$,\ $S_2E_1I_2$,\ $S_2E_1I_3$,\ $S_1S_2I_1$,\ $S_1S_2I_2$ \\
$S_2E_2$ & $S_2E_2I_1$,\ $S_2E_2I_2$,\ $S_2E_2I_3$,\ $S_2^{2}I_1$,\ $S_2^{2}I_2$,\ $S_2^{2}I_3$ \\
$S_2E_3$ & $S_2E_3I_1$,\ $S_2E_3I_2$,\ $S_2E_3I_3$,\ $S_2S_3I_2$,\ $S_2S_3I_3$ \\
$S_3E_2$ & $S_3E_2I_2$,\ $S_3E_2I_3$,\ $S_2S_3I_1$,\ $S_2S_3I_2$,\ $S_2S_3I_3$ \\
$S_3E_3$ & $S_3E_3I_2$,\ $S_3E_3I_3$,\ $S_3^{2}I_2$,\ $S_3^{2}I_3$ \\
\bottomrule
\end{longtable}


Across the 14 scalar derivatives there are \NumRawCubic{} monomial appearances.
Before the final seed set is formed, all products are written in a canonical
commutative form.
